\Needspace{8\baselineskip}\originalchapter{2006 作业九: Vlasov 与能量矩}
题源: 30 Nov 2006, due 8 Dec.
\begin{exercise}\label{ex:06-9-1}
原题1. 相空间粒子守恒 $D(f\Delta r\Delta v)/Dt=0$. (a)沿特征线差商 $f(r+v\Delta t,v+a\Delta t,t+\Delta t)$求Df/Dt; (b)变换r\ensuremath{{}^{\prime}}=r+v\ensuremath{\Delta}t,v\ensuremath{{}^{\prime}}=v+a(r,v,t)\ensuremath{\Delta}t, 用六维Jacobian求体积导数, 完成Vlasov推导.
\end{exercise}
\begin{solution}
(a)多元Taylor给 $Df/Dt=\partial_t f+\mathbf v\cdot\nabla_r f+\mathbf a\cdot\nabla_v f$. (b)变换矩阵为 $I+\Delta t A$, 行列式 $1+\Delta t\operatorname{tr}A+O(\Delta t^2)$, 其中 $\operatorname{tr}A=\nabla_r\cdot\mathbf v+\nabla_v\cdot\mathbf a$. r/v为独立坐标, $\nabla_r\cdot\mathbf v=0$; Lorentz 加速度 $\mathbf a=q(\mathbf E+\mathbf v\times\mathbf B)/m$ 在v空间的散度为零, 因反对称叉乘矩阵 $\operatorname{tr}A=0$. 故相空间体积保持, 守恒式除以体积得 $Df/Dt=0$. 非Lorentz阻力可有 $\nabla_v\cdot\mathbf a\ne0$, 此时不能沿轨道说f本身保持.
\end{solution}
\begin{exercise}\label{ex:06-9-2}
原题2. 从总能量矩
$\partial_t(mnu^2/2+3nT/2)+\nabla\cdot(mnu^2\mathbf u/2+5nT\mathbf u/2+\boldsymbol\Pi\cdot\mathbf u)+\nabla\cdot\mathbf q-q_s n\mathbf E\cdot\mathbf u=\int mv^2 C/2\dd^3v$,
用连续性、速度点乘动量和张量散度恒等式, 得 $3nDT/(2Dt)+p\nabla\cdot u=-\Pi_{ij}\partial_ju_i-\nabla\cdot q+Q$, $Q=\int mw^2C/2\dd^3v$.
\end{exercise}
\begin{solution}
连续性给 $\partial_t(nG)+\nabla\cdot(nG\mathbf u)=nDG/Dt$. 动量点乘u给 $mnD(u^2/2)/Dt=q_snE\cdot u-u\cdot\nabla p-u\cdot(\nabla\cdot\Pi)+u\cdot R$. 总能量减此动能, 利用\ensuremath{\nabla}·(pu)\ensuremath{-}u·\ensuremath{\nabla}p=p\ensuremath{\nabla}·u, 及 $\nabla\cdot(\Pi\cdot u)-u\cdot\nabla\cdot\Pi=\Pi_{ij}\partial_ju_i$, 得题式. 碰撞总能量减u·R再加/减连续性碰撞项, 正是以w=v\ensuremath{-}u写的Q; 磁力u·(u\ensuremath{\times}B)0不做功. \ensuremath{\Pi}此处为偏离标量p的黏性/各向异性压力张量, 与作业\ref{ex:03-5-2}的\textbf{总动量通量}同字母而含义不同.
\end{solution}
\begin{exercise}\label{ex:06-9-3}
原题3. 波 $\exp(-\ii\omega t+\ii kx)$、k实\ensuremath{\omega}可复, 电子归一一维Lorentz $F_e=v_e/[\pi(v_x^2+v_e^2)]$. (a)不动离子求代数色散; (b)解\ensuremath{\omega}(k); (c)可动离子 $F_i=v_i/[\pi(v_x^2+v_i^2)]$, |\ensuremath{\omega}/k|\ensuremath{\ll}ve求离子声波. 原件附上下半平面实轴积分公式.
\end{exercise}
\begin{solution}
按章\ref{ch:landau}由初值因果延拓, (a) $1-\omega_{pe}^2/(\omega+\ii|k|v_e)^2=0$. (b) $\omega=\pm\omega_{pe}-\ii|k|v_e$, 衰减支不改用另一个未经延拓的实轴积分. (c)电子近似屏蔽, $1+\omega_{pe}^2/(k^2v_e^2)-\omega_{pi}^2/(\omega+\ii|k|v_i)^2=0$, 得
\begin{equation}
 \omega=\pm\frac{\omega_{pi}}{\sqrt{1+\omega_{pe}^2/(k^2v_e^2)}}-\ii|k|v_i.
\end{equation}
长波实相速 $v_e\sqrt{m_e/m_i}$, 仍需满足原低频条件. Lorentz宽度不是有限温度, 上下半平面的普通实轴积分不能直接切换以绕过Landau轮廓.
\end{solution}
